When antitrust enforcement dismantles a procurement cartel, the convicted firms' workers disperse to competitors. Does this labor reallocation transmit collusive know-how to the receiving firms---or does it simply redistribute operational capacity? We test this by linking seven convicted cartels in S\~{a}o Paulo to 2.6~million public auctions and the universe of matched employer--employee records. Using a stacked difference-in-differences design, we find that firms absorbing ex-cartel workers increase their win rates by 5--6 percentage points in the same product markets---but do not raise their prices. An exit placebo confirms that the win-rate effect is cartel-specific: firms absorbing workers from non-cartel exits show a 3.7 times smaller effect. The price null, by contrast, is generic: any firm that absorbs workers from an exiting competitor charges more, regardless of cartel status---and the apparent price effect in a standard TWFE specification is an artifact of staggered-treatment bias. Characterizing the movers sharpens the mechanism: the workers who transfer are predominantly operational staff in service and sales roles, not the managers or professionals who might carry coordination strategies. The combined evidence points to capacity transfer, not collusion contagion. Enforcement dismantles the cartel and, through labor reallocation, strengthens the cartel's former rivals---without transmitting the ability to collude.
